%------------------------------------------------------------------------------
% 3.1 منهجية التطوير
%------------------------------------------------------------------------------
\section{منهجية التطوير}
\label{sec:methodology}

اعتمد فريق العمل منهجية التطوير الرشيقة \EN{Agile} \cite{schwaber2020scrum} في بناء نظام «بيان»، وذلك
لمواءمتها مع طبيعة المشروع الذي يجمع بين مكونات برمجية وتقنيات ذكاء اصطناعي
متقدمة، مثل توليد استعلامات قواعد البيانات من اللغة الطبيعية \EN{NL2SQL}،
والاسترجاع المعزز بالتوليد \EN{RAG}، ومعالجة اللغة العربية، وإدارة الأنظمة
متعددة الوكلاء المنسقة باستخدام إطار \EN{LangGraph}.

وتتطلب هذه التقنيات تكراراً مستمراً في تحسين النماذج اللغوية الكبيرة
\EN{LLMs}، والمطالبات \EN{Prompts}، وهيكلة المعرفة، فضلاً عن التحقق المتكرر
من دقة النتائج وزمن الاستجابة. لذلك، تقسم المنهجية الرشيقة نطاق العمل إلى
دورات تطوير قصيرة ومحددة زمنياً تسمى التكرارات أو الدورات \EN{Iterations / Sprints}.
وفي نهاية كل دورة، تُراجع المخرجات وتُقيّم قبل الانتقال إلى الدورة اللاحقة،
مما يتيح التكيف السريع مع التغيرات في المتطلبات أو مع نتائج تقييم النماذج
اللغوية دون التأثير الجوهري في الجدول الزمني العام للمشروع.

\subsection{دورة التطوير الرشيقة}
\label{subsec:agile-lifecycle}

تتكون كل دورة تطوير في مشروع «بيان» من ست مراحل متتابعة، مع إمكانية تكرارها
عند ظهور متطلبات جديدة أو الحاجة إلى تحسين أحد المكونات. وتضمن هذه المراحل
الانتقال المنظم من تحديد نطاق العمل إلى التقييم والتحسين، كما هو موضح في
الجدول \ref{tab:agile-phases}.

\begin{table}[H]
  \centering
  \caption{مراحل دورة التطوير الرشيقة المطبقة في مشروع «بيان».}
  \label{tab:agile-phases}
  \small
  \renewcommand{\arraystretch}{1.35}
  \begin{tabularx}{\textwidth}{>{\centering\arraybackslash}p{2.8cm}>{\raggedright\arraybackslash}X>{\raggedright\arraybackslash}p{4.1cm}}
    \toprule
    \textbf{المرحلة} & \textbf{أبرز الأنشطة} & \textbf{المخرج الرئيسي} \\
    \midrule
    \textbf{التخطيط}\\[-2pt]\EN{Planning}
    & تحديد أهداف الدورة الحالية، وإعداد قائمة مهام الدورة \EN{Sprint Backlog}،
      ومراجعة المتطلبات الوظيفية وغير الوظيفية المرتبطة بالمكون المستهدف.
    & وثيقة نطاق الدورة وقائمة المهام المحددة. \\
    \midrule
    \textbf{التصميم}\\[-2pt]\EN{Design}
    & نمذجة المعمارية الفنية للوكيل المستهدف، وتحديد هياكل البيانات والواجهات
      البرمجية \EN{APIs}، وإعداد التصميم الأولي لهندسة المطالبات \EN{Prompt Engineering}.
    & مخططات التدفق، وهياكل البيانات، ومخطط هندسة المطالبات. \\
    \midrule
    \textbf{التطوير}\\[-2pt]\EN{Development}
    & كتابة الشفرة البرمجية باستخدام \EN{Python} و\EN{FastAPI} و\EN{LangGraph}
      للواجهة الخلفية، أو \EN{Next.js} للواجهة الأمامية، وربط الوكلاء بمصادر
      البيانات \EN{PostgreSQL} و\EN{Qdrant}.
    & مكون برمجي جاهز يعمل بصورة مستقلة، مثل وكيل متخصص أو واجهة النظام. \\
    \midrule
    \textbf{الاختبار}\\[-2pt]\EN{Testing}
    & تنفيذ اختبارات الوحدة \EN{Unit Testing} واختبارات التكامل
      \EN{Integration Testing}، ومعايرة دقة الاستعلامات، والحد من الهلوسة،
      وقياس مؤشرات الأداء المعتمدة.
    & تقرير اختبارات موثق يوضح نسب الدقة وزمن الاستجابة. \\
    \midrule
    \textbf{الإطلاق والدمج}\\[-2pt]\EN{Deployment}
    & بناء المكونات البرمجية وحزمها داخل حاويات \EN{Docker Compose}، ودمجها
      في البيئة الموحدة، والتحقق من توافقها مع بقية الوكلاء.
    & نسخة محدثة من النظام، مدمجة ومجربة في بيئة التطوير. \\
    \midrule
    \textbf{المراجعة والتقييم}\\[-2pt]\EN{Review}
    & تقديم عرض توضيحي للمخرجات أمام مشرف المشروع، وتوثيق الدروس المستفادة،
      وتحديث الأولويات وقائمة مهام الدورة التالية.
    & قائمة تحسينات وتعديلات معتمدة للدورة المباشرة التالية. \\
    \bottomrule
  \end{tabularx}
\end{table}

\subsection{الدورات التنفيذية للمشروع}
\label{subsec:project-sprints}

لضمان التنفيذ المنهجي لمكونات نظام «بيان»، قُسّم نطاق المشروع إلى أربع دورات
تطوير تكرارية متتالية، بحيث تركز كل دورة على مكون رئيسي أو مجموعة مترابطة من
المكونات، وتنتهي بمخرج قابل للتقييم والدمج مع بقية النظام.

\begin{figure}[htbp]
  \centering
  \resizebox{0.91\textwidth}{!}{%==============================================================================
% Professional vertical roadmap for Bayan's four development sprints.
% Vector-only TikZ artwork; no external images or icon packages are required.
%==============================================================================

\begin{tikzpicture}[
  x=1cm,
  y=1cm,
  line cap=round,
  line join=round,
  connector/.style={->, >=stealth, line width=1.25pt, draw=kaudark!65},
  card title/.style={font=\Large\bfseries, align=center},
  card body/.style={font=\normalsize, text=kaudark, align=center},
  number/.style={font=\fontsize{19}{22}\selectfont\bfseries, text=white},
]

%------------------------------------------------------------------------------
% Shared card geometry
%------------------------------------------------------------------------------
\def\CardLeft{-7.35}
\def\CardRight{7.35}
\def\CardHalfHeight{1.23}
\def\IconX{-6.05}
\def\TextX{0.25}
\def\BadgeX{6.95}

%==============================================================================
% Sprint 1 — SQL Agent and Smart Views
%==============================================================================
\def\CardY{5.25}
\definecolor{sprintone}{HTML}{17639A}

% Soft shadow and card surface
\fill[black!10, rounded corners=8pt]
  (\CardLeft+0.10,\CardY-\CardHalfHeight-0.10)
  rectangle (\CardRight+0.10,\CardY+\CardHalfHeight-0.10);
\filldraw[fill=sprintone!4, draw=sprintone, line width=1.15pt, rounded corners=8pt]
  (\CardLeft,\CardY-\CardHalfHeight)
  rectangle (\CardRight,\CardY+\CardHalfHeight);

% Icon medallion
\fill[white] (\IconX,\CardY) circle (0.91);
\draw[sprintone!65, line width=1.15pt] (\IconX,\CardY) circle (0.91);
\draw[sprintone!18, line width=3.5pt] (\IconX,\CardY) circle (0.98);

% Database icon
\filldraw[fill=sprintone!82, draw=sprintone, line width=0.7pt]
  (\IconX-0.48,\CardY+0.36) arc[start angle=180,end angle=360,x radius=0.48,y radius=0.16]
  -- (\IconX+0.48,\CardY-0.29)
  arc[start angle=0,end angle=-180,x radius=0.48,y radius=0.16] -- cycle;
\draw[white!90, line width=0.65pt]
  (\IconX-0.48,\CardY+0.05) arc[start angle=180,end angle=360,x radius=0.48,y radius=0.16];
\draw[white!90, line width=0.65pt]
  (\IconX-0.48,\CardY-0.25) arc[start angle=180,end angle=360,x radius=0.48,y radius=0.16];
\filldraw[fill=white, draw=sprintone, line width=0.65pt, rounded corners=1.5pt]
  (\IconX+0.02,\CardY-0.52) rectangle (\IconX+0.63,\CardY-0.06);
\node[font=\scriptsize\bfseries, text=sprintone] at (\IconX+0.325,\CardY-0.29) {\EN{SQL}};

% Title, divider, and body
\node[card title, text=sprintone] at (\TextX,\CardY+0.52) {\RL{دورة التطوير الأولى}};
\draw[sprintone!70, line width=0.75pt] (-3.35,\CardY+0.15) -- (3.85,\CardY+0.15);
\fill[sprintone] (0.25,\CardY+0.15) circle (0.055);
\node[card body] at (\TextX,\CardY-0.32)
  {\RL{تطوير وكيل قواعد البيانات وطبقة الرؤى الذكية}\\[-1pt]
   {\small\color{sprintone}\EN{SQL Agent \& Smart Views}}};

% Number badge
\filldraw[fill=sprintone, draw=sprintone!75!black, line width=0.8pt]
  (\CardRight-0.08,\CardY+0.64) -- (\BadgeX-0.98,\CardY+0.64) --
  (\BadgeX-1.36,\CardY) -- (\BadgeX-0.98,\CardY-0.64) --
  (\CardRight-0.08,\CardY-0.64) -- cycle;
\node[number] at (\BadgeX-0.48,\CardY) {01};

%==============================================================================
% Sprint 2 — RAG Agent and document indexing
%==============================================================================
\def\CardY{1.75}
\definecolor{sprinttwo}{HTML}{2F7D32}

\fill[black!10, rounded corners=8pt]
  (\CardLeft+0.10,\CardY-\CardHalfHeight-0.10)
  rectangle (\CardRight+0.10,\CardY+\CardHalfHeight-0.10);
\filldraw[fill=sprinttwo!4, draw=sprinttwo, line width=1.15pt, rounded corners=8pt]
  (\CardLeft,\CardY-\CardHalfHeight)
  rectangle (\CardRight,\CardY+\CardHalfHeight);

\fill[white] (\IconX,\CardY) circle (0.91);
\draw[sprinttwo!65, line width=1.15pt] (\IconX,\CardY) circle (0.91);
\draw[sprinttwo!18, line width=3.5pt] (\IconX,\CardY) circle (0.98);

% Document and search icon
\filldraw[fill=sprinttwo!88, draw=sprinttwo, line width=0.7pt, rounded corners=1.5pt]
  (\IconX-0.46,\CardY-0.55) rectangle (\IconX+0.28,\CardY+0.54);
\fill[white] (\IconX+0.02,\CardY+0.54) -- (\IconX+0.28,\CardY+0.28) --
  (\IconX+0.02,\CardY+0.28) -- cycle;
\draw[white, line width=0.75pt]
  (\IconX-0.29,\CardY+0.08) -- (\IconX+0.08,\CardY+0.08)
  (\IconX-0.29,\CardY-0.12) -- (\IconX+0.03,\CardY-0.12)
  (\IconX-0.29,\CardY-0.32) -- (\IconX-0.03,\CardY-0.32);
\draw[sprinttwo, line width=1.1pt] (\IconX+0.34,\CardY-0.21) circle (0.27);
\draw[sprinttwo, line width=1.4pt] (\IconX+0.54,\CardY-0.41) -- (\IconX+0.75,\CardY-0.62);

\node[card title, text=sprinttwo] at (\TextX,\CardY+0.52) {\RL{دورة التطوير الثانية}};
\draw[sprinttwo!70, line width=0.75pt] (-3.35,\CardY+0.15) -- (3.85,\CardY+0.15);
\fill[sprinttwo] (0.25,\CardY+0.15) circle (0.055);
\node[card body] at (\TextX,\CardY-0.32)
  {\RL{تطوير وكيل البحث وفهرسة المستندات واللوائح}\\[-1pt]
   {\small\color{sprinttwo}\EN{RAG Agent \& Document Indexing}}};

\filldraw[fill=sprinttwo, draw=sprinttwo!75!black, line width=0.8pt]
  (\CardRight-0.08,\CardY+0.64) -- (\BadgeX-0.98,\CardY+0.64) --
  (\BadgeX-1.36,\CardY) -- (\BadgeX-0.98,\CardY-0.64) --
  (\CardRight-0.08,\CardY-0.64) -- cycle;
\node[number] at (\BadgeX-0.48,\CardY) {02};

%==============================================================================
% Sprint 3 — Multi-agent integration and Arabic UI
%==============================================================================
\def\CardY{-1.75}
\definecolor{sprintthree}{HTML}{633A98}

\fill[black!10, rounded corners=8pt]
  (\CardLeft+0.10,\CardY-\CardHalfHeight-0.10)
  rectangle (\CardRight+0.10,\CardY+\CardHalfHeight-0.10);
\filldraw[fill=sprintthree!4, draw=sprintthree, line width=1.15pt, rounded corners=8pt]
  (\CardLeft,\CardY-\CardHalfHeight)
  rectangle (\CardRight,\CardY+\CardHalfHeight);

\fill[white] (\IconX,\CardY) circle (0.91);
\draw[sprintthree!65, line width=1.15pt] (\IconX,\CardY) circle (0.91);
\draw[sprintthree!18, line width=3.5pt] (\IconX,\CardY) circle (0.98);

% Supervisor/network icon
\foreach \ang in {30,90,150,210,270,330} {
  \draw[sprintthree!75, line width=0.8pt]
    (\IconX,\CardY) -- ++(\ang:0.56);
  \filldraw[fill=white, draw=sprintthree, line width=0.9pt]
    (\IconX,\CardY) ++(\ang:0.62) circle (0.12);
}
\filldraw[fill=sprintthree, draw=sprintthree!80!black, line width=0.7pt]
  (\IconX,\CardY) circle (0.34);
\fill[white] (\IconX,\CardY+0.10) circle (0.09);
\fill[white] (\IconX-0.16,\CardY-0.16)
  .. controls (\IconX-0.10,\CardY-0.03) and (\IconX+0.10,\CardY-0.03) ..
  (\IconX+0.16,\CardY-0.16) -- cycle;

\node[card title, text=sprintthree] at (\TextX,\CardY+0.52) {\RL{دورة التطوير الثالثة}};
\draw[sprintthree!70, line width=0.75pt] (-3.35,\CardY+0.15) -- (3.85,\CardY+0.15);
\fill[sprintthree] (0.25,\CardY+0.15) circle (0.055);
\node[card body] at (\TextX,\CardY-0.32)
  {\RL{دمج الوكلاء وبناء الواجهة العربية التفاعلية}\\[-1pt]
   {\small\color{sprintthree}\EN{Supervisor Agent \& Arabic Web Interface}}};

\filldraw[fill=sprintthree, draw=sprintthree!75!black, line width=0.8pt]
  (\CardRight-0.08,\CardY+0.64) -- (\BadgeX-0.98,\CardY+0.64) --
  (\BadgeX-1.36,\CardY) -- (\BadgeX-0.98,\CardY-0.64) --
  (\CardRight-0.08,\CardY-0.64) -- cycle;
\node[number] at (\BadgeX-0.48,\CardY) {03};

%==============================================================================
% Sprint 4 — Testing, optimization, and deployment
%==============================================================================
\def\CardY{-5.25}
\definecolor{sprintfour}{HTML}{D66A00}

\fill[black!10, rounded corners=8pt]
  (\CardLeft+0.10,\CardY-\CardHalfHeight-0.10)
  rectangle (\CardRight+0.10,\CardY+\CardHalfHeight-0.10);
\filldraw[fill=sprintfour!4, draw=sprintfour, line width=1.15pt, rounded corners=8pt]
  (\CardLeft,\CardY-\CardHalfHeight)
  rectangle (\CardRight,\CardY+\CardHalfHeight);

\fill[white] (\IconX,\CardY) circle (0.91);
\draw[sprintfour!65, line width=1.15pt] (\IconX,\CardY) circle (0.91);
\draw[sprintfour!18, line width=3.5pt] (\IconX,\CardY) circle (0.98);

% Gear/check icon
\foreach \ang in {0,45,...,315} {
  \fill[sprintfour] (\IconX,\CardY) ++(\ang:0.48)
    ++(-0.11,-0.13) rectangle ++(0.22,0.26);
}
\filldraw[fill=sprintfour, draw=sprintfour!80!black, line width=0.7pt]
  (\IconX,\CardY) circle (0.49);
\fill[white] (\IconX,\CardY) circle (0.32);
\draw[sprintfour, line width=1.35pt]
  (\IconX-0.18,\CardY) -- (\IconX-0.04,\CardY-0.14) --
  (\IconX+0.23,\CardY+0.17);

\node[card title, text=sprintfour] at (\TextX,\CardY+0.52) {\RL{دورة التطوير الرابعة}};
\draw[sprintfour!70, line width=0.75pt] (-3.35,\CardY+0.15) -- (3.85,\CardY+0.15);
\fill[sprintfour] (0.25,\CardY+0.15) circle (0.055);
\node[card body] at (\TextX,\CardY-0.32)
  {\RL{الاختبار الشامل وضبط الأداء والتجهيز للنشر}\\[-1pt]
   {\small\color{sprintfour}\EN{Testing, Optimization \& Deployment}}};

\filldraw[fill=sprintfour, draw=sprintfour!75!black, line width=0.8pt]
  (\CardRight-0.08,\CardY+0.64) -- (\BadgeX-0.98,\CardY+0.64) --
  (\BadgeX-1.36,\CardY) -- (\BadgeX-0.98,\CardY-0.64) --
  (\CardRight-0.08,\CardY-0.64) -- cycle;
\node[number] at (\BadgeX-0.48,\CardY) {04};

%------------------------------------------------------------------------------
% Connectors — drawn last to remain crisp between cards
%------------------------------------------------------------------------------
\draw[connector] (0,4.00) -- (0,3.03);
\draw[connector] (0,0.50) -- (0,-0.47);
\draw[connector] (0,-3.00) -- (0,-3.97);

% Subtle decorative dot fields at the outer margins
\foreach \y/\c in {5.25/sprintone,1.75/sprinttwo,-1.75/sprintthree,-5.25/sprintfour} {
  \foreach \dx in {0,0.18,0.36} {
    \foreach \dy in {-0.18,0,0.18} {
      \fill[\c!48] (-7.82+\dx,\y+\dy) circle (0.025);
      \fill[\c!48] (7.82-\dx,\y+\dy) circle (0.025);
    }
  }
}

\end{tikzpicture}
}
  \caption{خارطة الدورات التطويرية الأربع المعتمدة في مشروع «بيان».}
  \label{fig:project-sprints}
\end{figure}

\subsubsection{الدورة الأولى: بناء وكيل استعلام قواعد البيانات}
\label{subsubsec:sprint-one}

\textbf{المستهدف:} إعداد المعمارية الأساسية للتعامل مع البيانات المهيكلة في
قاعدة بيانات \EN{PostgreSQL}، مع ضمان أمن الاستعلامات.

\textbf{الأنشطة:} تطوير طبقة الرؤى الذكية \EN{Smart Views} لعزل الجداول
المعقدة، وبناء محرك تحويل النص إلى استعلام \EN{NL2SQL Engine}، وإدراج آليات
العزل للقراءة فقط \EN{Read-Only Sandboxing} لمنع التعديل غير المصرح به، ثم
اختبار دقة التوليد على استعلامات مالية وإدارية وأكاديمية.

\textbf{المخرج:} وكيل \EN{SQL} مستقل يعمل عبر \EN{FastAPI} ويسترجع البيانات
بدقة من الجداول المحددة مع الالتزام بقيود الوصول الآمن.

\subsubsection{الدورة الثانية: بناء وكيل البحث في المستندات واللوائح}
\label{subsubsec:sprint-two}

\textbf{المستهدف:} معالجة البيانات غير المهيكلة، بما في ذلك اللوائح والسياسات
والمستندات التنظيمية الخاصة بالجامعة والمؤسسات ذات الصلة.

\textbf{الأنشطة:} إعداد قاعدة البيانات المتجهة \EN{Qdrant Vector DB}، وتطوير
معالج للنصوص العربية وتجزئتها \EN{Text Chunking}، وإنشاء المتجهات
\EN{Embeddings} باستخدام نماذج داعمة للغة العربية، وبناء آلية استرجاع معزز
بالمصادر ترفق اسم المستند ورقم الصفحة مع الإجابة متى توفرت البيانات.

\textbf{المخرج:} وكيل \EN{RAG} متكامل يجيب عن الأسئلة النصية بإجابات دقيقة
ومدعومة بمراجع واضحة.

\subsubsection{الدورة الثالثة: دمج الوكلاء وبناء الواجهة الأمامية}
\label{subsubsec:sprint-three}

\textbf{المستهدف:} تحقيق التكامل بين المكونات البرمجية وإتاحة واجهة مستخدم
عربية سلسة.

\textbf{الأنشطة:} بناء الوكيل المشرف \EN{Supervisor Agent} باستخدام إطار
\EN{LangGraph} لتوجيه الاستفسارات \EN{Intent Routing} إلى الوكيل المختص
(وكيل \EN{SQL} أو وكيل \EN{RAG} أو كليهما)، وإدارة ذاكرة السياق
\EN{Context Management}، ثم تطوير واجهة محادثة تفاعلية باللغة العربية
وبدعم الاتجاه من اليمين إلى اليسار \EN{RTL} باستخدام \EN{Next.js} و\EN{Tailwind CSS}.

\textbf{المخرج:} نظام موحد متعدد الوكلاء يعمل عبر واجهة ويب عربية حوارية
متكاملة.

\subsubsection{الدورة الرابعة: الاختبار الشامل وضبط الأداء والنشر}
\label{subsubsec:sprint-four}

\textbf{المستهدف:} التحقق من موثوقية النظام، وإدارة المخاطر التقنية، وتجهيز
البيئة التشغيلية للنشر والتقييم.

\textbf{الأنشطة:} تنفيذ اختبارات شاملة من البداية إلى النهاية
\EN{End-to-End Testing}، وتقييم زمن الاستجابة وجودة الإجابات، وضبط أداء
النموذج اللغوي المحلي \EN{Qwen 2.5 14B} عبر \EN{Ollama}، ثم حزم النظام
بالكامل باستخدام \EN{Docker} و\EN{Docker Compose}.

\textbf{المخرج:} بيئة تشغيل موحدة قابلة للنشر السريع، تحتوي على مكونات نظام
«بيان» في حاويات مستقلة وجاهزة للتسليم والتطبيق التجريبي.

\subsection{ملخص الدورات ومخرجاتها}
\label{subsec:sprints-summary}

يلخص الجدول \ref{tab:sprints-summary} العلاقة بين كل دورة تطوير والمكون
المستهدف والأدوات المستخدمة والمخرج المباشر المتحقق منها.

\begin{table}[H]
  \centering
  \caption{الملخص التنفيذي للدورات الأربع في منهجية تطوير نظام «بيان».}
  \label{tab:sprints-summary}
  \footnotesize
  \renewcommand{\arraystretch}{1.55}
  \begin{tabularx}{\textwidth}{>{\centering\arraybackslash}p{1.5cm}>{\raggedright\arraybackslash}p{3.2cm}>{\raggedright\arraybackslash}X>{\raggedright\arraybackslash}p{3.8cm}}
    \toprule
    \textbf{الدورة} & \textbf{المكون المستهدف} & \textbf{الأنشطة والأدوات الرئيسية} & \textbf{المخرج النهائي المباشر} \\
    \midrule
    \EN{Sprint 1} & وكيل \EN{SQL} وطبقة \EN{Smart Views}
    & \EN{PostgreSQL}، و\EN{FastAPI}، ومحرك \EN{NL2SQL}، وطبقة الرؤى الذكية.
    & وكيل \EN{SQL} آمن ومستقل لاسترجاع البيانات المهيكلة. \\
    \midrule
    \EN{Sprint 2} & وكيل \EN{RAG} والبحث في المستندات
    & قاعدة \EN{Qdrant} المتجهة، والمتجهات الدلالية العربية، وتجزئة المستندات.
    & وكيل \EN{RAG} يجيب استناداً إلى السياسات واللوائح مع توثيق المصدر. \\
    \midrule
    \EN{Sprint 3} & دمج الوكلاء والواجهة الأمامية
    & الوكيل المشرف \EN{LangGraph}، و\EN{Next.js}، و\EN{Tailwind CSS}.
    & نظام هجين موحد يعمل عبر واجهة ويب عربية حوارية. \\
    \midrule
    \EN{Sprint 4} & الاختبار والتحسين والنشر
    & \EN{Docker Compose}، و\EN{Ollama (Qwen 2.5 14B)}، واختبارات
      \EN{End-to-End}.
    & نسخة تشغيلية محزمة وجاهزة للتقييم والتطبيق التجريبي. \\
    \bottomrule
  \end{tabularx}
\end{table}

وبذلك تحقق المنهجية تدرجاً عملياً يبدأ ببناء الوكلاء المتخصصين، ثم دمجهم في
معمارية متعددة الوكلاء، وينتهي بالاختبار الشامل والتجهيز للنشر. كما تسمح
المخرجات المرحلية بتقييم كل مكون بصورة مستقلة قبل اعتماده ضمن النظام الموحد.
